\section{Nanbeige4.3-33B公開、エージェント性能に関心}
Nanbeige LLM Labが、33Bパラメータの密モデル「Nanbeige4.3-33B」をApache 2.0でHugging Faceに公開した。エージェント向け教師モデルとして位置づけられ、モデルカードではTerminal-Bench 2.1のスコア80.9を自己報告し、Qwen3.8-27Bの73.0を上回ると主張する。X上ではローカルLLM利用者が性能やデプロイの現実性を議論した。ただし数値の独立再現は未確認であり、長いコンテキストでのKVキャッシュ消費も論点となっている。
\dailyxsource{Hikari（ローカルLLM関連投稿）}{https://x.com/Hikari_07_jp/status/2114350960730644847}

\section{Claude Duet、プロアクティブ機能の利用報告}
Anthropicの「Claude Duet」を巡り、ウェイトリスト通過や使用体験を伝える投稿が続いた。朝のブリーフやバックグラウンド作業を支援するエージェントアプリで、Grokの記録では米国のPro／Maxユーザー向け研究プレビューとして案内されている。外部操作には承認を必要とする仕組みが説明されている。今後の焦点は一般提供の時期と利用上の制限、既存エージェントとの実務的な違いにある。引用投稿は観測窓の終端に位置する。
\dailyxsource{Ramesh Donthaによる投稿}{https://x.com/EntrepreneursAI/status/2114477160140280108}

\section{AnthropicのIPOとモデル発表時期を巡る推測}
Anthropicの新規株式公開（IPO）の目標、収益規模、モデル発表のタイミングを結びつけた推測がXで共有された。投稿者はBloombergなどの報道を参照したとするが、Grokの収集では観測窓内の公式な裏付けや一次資料を確認できていない。成長率や収益の具体的な数字を確定情報として扱うことはできず、正式な申請書類や同社の発表を待つ必要がある。記録された一次X投稿時刻は07:04 JST頃で、観測窓外に当たる。
\dailyxsource{Mega Evolved Frogによる投稿（観測窓外）}{https://x.com/MegaEvolvedFrog/status/2114478290203238868}

\section{GPT-6.1 Astraの遅延を示唆する未確認投稿}
「GPT-6.1 Astraは今週ローンチされない」とする単発のX投稿が観測された。投稿には「Tiboが確認した」との表現や安全性に関する問題の主張があるものの、Grokの収集ではOpenAIによる公式言及や一次資料を確認できていない。したがって、提供時期の変更や安全性上の問題が事実として判明したわけではない。モデルのリリース予定については、今後の公式発表による確認が必要である。
\dailyxsource{Impact Factory Aiによる未確認の主張}{https://x.com/impactfactoryai/status/2114474432609616376}

\section{宮本茂氏の生成AIとゲーム制作に関する発言が話題}
任天堂の宮本茂氏が、生成AIによってデザイナー以外もゲームを制作できる時代を肯定的に捉えたとする引用がX上で拡散した。ゲーム制作に必要な技能や、創作への参加の広がりを巡る反応が見られる。ただしGrokは元インタビューの一次確認が限定的だと注記しており、発言の正確な文脈は未確認である。引用投稿の時刻は07:03 JST頃で、今回の観測窓をわずかに越えている。
\dailyxsource{Andrew Curranによる引用投稿（観測窓外）}{https://x.com/AndrewCurran_/status/2114478047030366561}

\section{10月のモデルリリースを振り返る投稿}
10月の生成AIモデルを総括するダイジェスト投稿が複数共有された。GPT-6、Claude Haiku 5.5、Gemini 4 Argonやオープンモデルなどを列挙し、モデル選択の幅が広がっていることを示す内容である。ただし、この日の新規発表をまとめたものではなく、既存リリースの振り返りが中心だった。列挙された各モデルの公開条件や性能についても、Grokは個別の一次確認が必要だとしている。速報と月次の総括を混同しないことが重要となる。
\dailyxsource{PriyaによるAIダイジェスト}{https://x.com/glitch_catAI/status/2114442370917744710}

\section{AI生成広告のラベリング法への再言及}
米ペンシルバニア州で、AIが生成・改変した広告の表示を義務づける法律が制定されたことが、X上のダイジェストで改めて取り上げられた。Grokは法律への署名自体は今回の観測窓より前だと明記しており、本号で扱うのは窓内での政策議論である。広告の透明性を求める動きとして注目される一方、施行時期や具体的な運用、他州への波及については今後の確認事項となる。
\dailyxsource{RC | Macro Siftによる言及（制定は観測窓外）}{https://x.com/themacrosift/status/2114400152954048615}

\section{オープンウェイトモデルの実行コストを議論}
オープンウェイトのモデルでも、規模が大きければ一般の利用者には実行が難しいという指摘がX上で共有された。モデルの重みが公開されていることと、手元の計算資源で実用的に動かせることは別の問題であり、量子化や蒸留による小型化の重要性が話題となった。Grokはこの投稿を技術的な意見として記録しており、具体的な性能比較の一次資料は確認していない。投稿時刻は07:06 JST頃で観測窓外である。
\dailyxsource{Eric Pereiraによる意見（観測窓外）}{https://x.com/iamericpereira/status/2114478680168701971}
